%% kapa: exam=1 round=20 year=2009 period=1 form=A subject=civil_law subject_ko=민법 seq=29 number=29
%% source: Q-Net 감정평가사 기출문제 (cst00309) artlSeq=5024178 file=extracted/1008/5024178/2009년 제20회 감정평가사 1차시험 기출문제/1교시.hwp
%% conversion: kordoc 4.24.51 -> tools/gen1.py (auto); review=proofread-sonnet; flags=none
\begin{question}{29}
다음 중 대리권 등에 관한 설명으로 옳지 \uline{않은} 것은?
\begin{choices}
  \item 대리인의 행위가 자기계약 및 쌍방대리금지 조항에 위배된 경우 무권대리행위로 된다.
  \item 대리권이 존재하는 것은 분명하지만, 그 범위가 불분명한 경우, 대리인은 보존행위만을 할 수 있다.
  \item 원인이 되는 기초적 법률관계가 종료되기 전이라도 본인은 언제든지 수권행위를 철회할 수 있다.
  \item 대리인이 사기나 강박을 당하지 않는 한, 본인이 사기나 강박을 당했더라도 본인은 대리행위를 취소할 수 없다.
  \item 금치산자라도 임의대리인이 될 수 있다.
\end{choices}
\end{question}
\answer{2}
